\section{VALIDATION PROTOCOLS}
\label{app:validation}

The analyses of \cref{sec:allegiance} share one design: a target with a known allegiance, probed identically from $\vmu_z$ and $\vmu_w$ (posterior means of a converged model), reported per type and pooled as cell-weighted means, with reference rows and spatial-block cross-validation \citep{roberts2017} using the training tiles as fold groups (five folds of contiguous tiles). Random cross-validation is invalid for every spatial target here: autocorrelation leaks the answer through neighbouring cells landing on both sides of the split. Every analysis is cheap post hoc and is run at every point of the leakage sweep with seed envelopes.

\paragraph{Reference rows.} (i) \emph{floor}: the target permuted within type; (ii) \emph{depth baseline}: the same probe from $\log\ell_i$ alone, because depth varies across tissue and a latent that encodes depth buys spatial skill technically (G2/M cells also carry about twice the mRNA \needsource, so depth is informative even for intrinsic targets); (iii) for expression-derived targets, a \emph{linear reference} from the leading $50$ principal components of the log-normalised counts. The linear reference is not a ceiling: a nonlinear encoder can and does exceed it. \todo{name it consistently as ``linear reference'' everywhere}

\paragraph{Cell-cycle phase (intrinsic).} Continuous S and G2/M scores from standard gene sets \citep[$18$ S and $34$ G2/M markers on the panel; split-half reliability $0.22$ and $0.52$]{tirosh2016}; hard phase calls default to G1 at this depth and are not used. Within-type ridge regression from each latent on the four types with the highest fraction of proliferation-marker-positive cells, reported as a per-type-centred cell-weighted pooled $R^2$; the mean over types is dominated by types whose within-type score variance is near zero and is reported only in the by-type table. \todo{consider gating types on within-type score variance}

\paragraph{Discrete niche (spatial).} Niche labels from $k$-means on neighbour composition over connected cells ($10$ niches; $6$ and $15$ as robustness), never from the latents. Types with at least two niches of at least $500$ cells are eligible; per type, a class-balanced multinomial logistic regression from each latent, reported as macro one-versus-rest area under the curve \citep{hand2001} and balanced accuracy. Expectation: $\vw$ well above the depth baseline; $\vz$ within noise of the larger of floor and depth baseline. The observed excess of $\vz$ over the depth baseline is reported as it is, and triaged per dimension by the Moran's I analysis below.

\paragraph{Moran's I per dimension (both latents).} Moran's I \citep{moran1950} of each dimension centred per type, so that the legitimate type mosaic is removed from both latents while the full graph is kept; the model's own pruned graph with binary row-normalised weights; a permutation null from at least $1{,}000$ shuffles of the centred values; dimensions whose variance is below $10^{-3}$ of the largest are flagged as collapsed and excluded from the mean. A high-autocorrelation $\vz$ dimension is a flag, not a failure: intrinsically spatial biology (clone patches, proliferative niches) is allowed, since $\vz$ claims independence of context-induced variance, not spatial uniformity. Each high-autocorrelation $\vz$ dimension is triaged by its in-sample $R^2$ on composition (leakage) and its correlation with the cycle scores (benign).

\paragraph{Pseudotime and the tissue gradient.} A principal curve \citep{hastie1989} in each latent space; its agreement with the tissue gradient measured as the $R^2$ of the pseudotime on neighbour composition and as the graph-neighbour coherence of the pseudotime, both raw and after partialling out per-type means. The raw values are dominated by type homophily and are reported only to show it.

\paragraph{Distance to landmarks (spatial, geometry-derived).} Landmark classes built from types and positions only: vasculature (endothelial cells, with pericytes included when they lie within $30\um$ of endothelium, clustered into instances by density-based clustering \citep{ester1996}), compact smooth-muscle structures (instances of $20$ to $500$ cells), and the tumour--stroma interface (from a $k$-nearest-neighbour-smoothed tumour fraction, taking the boundary). Target $\log(1 + \min(d, 500\um))$; the landmark's own constituent type excluded from its regression; bands \emph{near} ($< 30\um$, one hop, a composition echo and sanity check only), \emph{mid} ($50$ to $300\um$, the evidence, since one-hop context cannot see the landmark), and \emph{far} ($300$ to $500\um$, expected failure). Ridge regression fitted and scored within each band. Because the landmark sets are type-defined and types are expression-derived, a fourth reference row, the probe from raw neighbour composition, measures the definitional share of any signal. The probe direction per landmark is unit-normalised and compared across landmarks by cosine within one run, never across runs; gene signatures are read as $\vB\hat{\boldsymbol\theta}$ restricted to genes expressed in at least $1\%$ of cells. \todo{the interface landmark is roughly $60\%$ definitional and vasculature is a null on the existing run; decide whether the analysis appears beyond the appendix}

\paragraph{Allegiance matrix.} Rows: cell-cycle score, niche label, mid-band landmark distance, type-partialled pseudotime gradient. Columns: $\vmu_z$, $\vmu_w$, with floor, depth baseline and linear reference marked per cell. Expected pattern: block-diagonal.
